\clearpage
\section{Deferred Technical Details}
\label{appx:tech}

\subsection{Pseudocode}
\label{appx:pseudocode}

We give the protocols in order of composition:
Algorithms~\ref{alg:ip}, \ref{alg:hadamard}, and
\ref{alg:permutation} implement the packed reductions of
Section~\ref{sec:blocks}; Algorithm~\ref{alg:packed-pcs}
authenticates their terminal evaluations as in Section~\ref{sec:pcs};
and Algorithms~\ref{alg:single-instance-packing} and
\ref{alg:multiple-instance-packing} compose these components
for the two settings of Section~\ref{sec:together}.
For the intermediate-field variant of Section~\ref{sec:hadamard},
Algorithm~\ref{alg:hadamard} uses the Lagrange basis over
\(\mathbb K\) and computes its quotient via
Eq.~\eqref{eq:hadamard-small-field-aggregation}.
Tabulation changes this computation, not the messages or checks.

% !TEX root = ../main.tex
\begin{algorithm}[!t]
\caption{Packed inner-product proof}
\label{alg:ip}
\KwIn{
\(N=m\kappa\), where \(m, \kappa \in 2^\mathbb{N}\);
\(c\in\Field\);
\(\bm a,\bm b\in\Field^N\)
}
\KwOut{
Terminal values \(\ext{v}_1,\ext{v}_2\) and MLE points
\(\ext{\bm y}_a, \ext{\bm y}_b \in\ExtensionField^{\log_2 N}\);
acceptance bit \(\mathsf{ok}\)
}

Arrange \(\bm a,\bm b\) columnwise as
\(\bm A\leftarrow\concat{\bm a^{(1)},\ldots,\bm a^{(\kappa)}}\)
and \(\bm B\leftarrow\concat{\bm b^{(1)},\ldots,\bm b^{(\kappa)}}
\in\Field^{m\times\kappa}\)\;

\tcp{Compute the auxiliary coefficients over the base field}
$\Prover$ sets:\\
\Indp
\(
    \bm g_A(X)
    \leftarrow
    \sum_{j=1}^{\kappa}
    \bm a^{(j)}X^{j-1}
\);~
\(
    \bm g_B(X)
    \leftarrow
    \sum_{j=1}^{\kappa}
    \bm b^{(j)}X^{j-1}
\)\;

\ForEach{
\(\tau\in
\{-(\kappa-1),\ldots,-1,1,\ldots,\kappa-1\}\)
}{
    Set \(
        f_\tau
        \leftarrow
        \sum_{\substack{j,j'\in[\kappa]\\j-j'=\tau}}
        \left\langle
            \bm a^{(j)},\bm b^{(j')}
        \right\rangle
    \)\;
}

Set
\(
    h_+(X)
    \leftarrow
    \sum_{\tau=1}^{\kappa-1}
    f_\tau X^{\tau-1}
\);~
\(
    h_-(X)
    \leftarrow
    \sum_{\tau=1}^{\kappa-1}
    f_{-\tau}X^{\tau-1}
\)\;
\Indm

\(\PtoV\):
the coefficients of \(h_+(X)\) and \(h_-(X)\)\;

\(\VtoP\):
\(\ext{x}\ChooseRandom\ExtensionField^\times\)\;

\tcp{Reduce to a length-\(m\) extension-field inner product}
$\Prover$ sets
\(
    \ext{\bm a}
    \leftarrow
    \bm g_A(\ext{x})
    \in\ExtensionField^m
\),
\(
    \ext{\bm b}
    \leftarrow
    \bm g_B(\ext{x}^{-1})
    \in\ExtensionField^m
\)\;

$\Verifier$ sets
\(
    \ext{c}
    \leftarrow
    c
    +\ext{x}h_+(\ext{x})
    +\ext{x}^{-1}h_-(\ext{x}^{-1})
\)\;

$\Prover$ and $\Verifier$ apply Fact~\ref{fact:ip} to
\(
    \ext{c}
    \EqOrNot
    \left\langle
        \ext{\bm a},\ext{\bm b}
    \right\rangle,
\)
and obtain terminal values
\(
    \ext{v}_1,\ext{v}_2
\),
the challenge point
\(
    \ext{\bm y}'\in
    \ExtensionField^{\log_2m}
\),
and the acceptance bit \(\mathsf{ok}\)\;

\(\Prover\) and \(\Verifier\) set:\\
\Indp
\(
    \ext{\bm y}_a
    \leftarrow
    \bigl(
        \ext{x}^{2^{\log_2\kappa-1}},
        \ext{x}^{2^{\log_2\kappa-2}},
        \ldots,
        \ext{x}
    \bigr)
    \mathbin{\|}
    \ext{\bm y}' \in \ExtensionField^{\log_2 N}
\)\;
\(
    \ext{\bm y}_b
    \leftarrow
    \bigl(
        \ext{x}^{-2^{\log_2\kappa-1}},
        \ext{x}^{-2^{\log_2\kappa-2}},
        \ldots,
        \ext{x}^{-1}
    \bigr)
    \mathbin{\|}
    \ext{\bm y}' \in \ExtensionField^{\log_2 N}
\)\;
\Indm

Output the terminal claims
\(
    \ext{v}_1
    \EqOrNot
    \left\langle
        \bm a,
        \tensorspan{\ext{\bm y}_a}
    \right\rangle
\) and
\(
    \ext{v}_2
    \EqOrNot
    \left\langle
        \bm b,
        \tensorspan{\ext{\bm y}_b}
    \right\rangle
\)\;

\Return
\(
    (\mathsf{ok},\ext{v}_1,\ext{v}_2,\ext{\bm y}_a, \ext{\bm y}_b)
\)\;
\end{algorithm}

% !TEX root = ../main.tex
\begin{algorithm}[!t]
\caption{Packed Hadamard-product proof}
\label{alg:hadamard}
\KwIn{
\(N=m\kappa\), \(m\in2^{\mathbb N}\), \(\kappa\geq2\);
a fixed claim \(\ext c\in\ExtensionField\) and matrices
\(\bm A,\bm B\in\Field^{m\times\kappa}\) fixed by the input commitments;
column weights \(\ext{\bm r}\in\ExtensionField^\kappa\);
row weights
\(\ext{\bm l}=\bigotimes_{j=1}^{\log_2m}
(\ext l_{j,0},\ext l_{j,1})^\top\);
nodes \(\mathcal H=\{\gamma_1,\ldots,\gamma_\kappa\}\subseteq\mathbb K\),
with \(\Field\subseteq\mathbb K\subseteq\ExtensionField\) and
\(|\mathbb K|\geq2\kappa-1\)
}
\KwOut{
Terminal values \(\ext v_1,\ext v_2\);
column weights \(\ext{\bm\rho}\in\ExtensionField^\kappa\);
tensor row weights \(\ext{\bm\xi}_1,\ext{\bm\xi}_2\in\ExtensionField^m\);
acceptance bit \(\mathsf{ok}\)
}

\tcp{Encode the weighted column claims}
\(\Prover\) computes
\(\ext{\bm c}^{\top}\leftarrow
\ext{\bm l}^{\top}(\bm A\circ\bm B)\)\;
In the Lagrange basis on \(\mathcal H\), define
\(\bm g_A(X)\leftarrow\sum_j\bm a^{(j)}\psi_j(X)\),
\(\bm g_B(X)\leftarrow\sum_j\bm b^{(j)}\psi_j(X)\), and
\(g_C(X)\leftarrow\sum_j\ext c_j\psi_j(X)\)\;
\(\Prover\) computes the residual
\(f(X)\leftarrow
\langle\ext{\bm l},\bm g_A(X)\circ\bm g_B(X)\rangle-g_C(X)\)
at \(\kappa-1\) fixed nodes in \(\mathbb K\setminus\mathcal H\),
interpolates with the zeros on \(\mathcal H\), and sets
\(h(X)\leftarrow f(X)/\Psi(X)\)\;
\(\PtoV\): \(\ext{\bm c}\) and the \(\kappa-1\) coefficients of \(h\)\;
\(\Verifier\) sets
\(\mathsf{ok}_0\leftarrow
(\ext c\EqOrNot\langle\ext{\bm c},\ext{\bm r}\rangle)\)\;
\(\VtoP\): \(\ext x\ChooseRandom\ExtensionField\setminus\mathcal H\)\;

\tcp{Reduce to a length-\(m\) extension-field inner product}
\(\Prover\) and \(\Verifier\) compute
\(\ext{\bm\rho}\leftarrow
(\psi_1(\ext x),\ldots,\psi_\kappa(\ext x))^\top\)\;
\(\Prover\) computes
\(\ext{\bm a}\leftarrow\bm A\ext{\bm\rho}\) and
\(\ext{\bm b}\leftarrow\bm B\ext{\bm\rho}\)\;
\(\Verifier\) computes
\(\ext c_x\leftarrow
\langle\ext{\bm c},\ext{\bm\rho}\rangle=g_C(\ext x)\)\;
\(\Prover\) and \(\Verifier\) apply Fact~\ref{fact:ip} to
\(\langle\ext{\bm a},\ext{\bm b}\circ\ext{\bm l}\rangle
\EqOrNot\ext c_x+\Psi(\ext x)h(\ext x)\),
obtaining \(\mathsf{ok}_{\mathsf{IP}}\),
\(\ext v_1,\ext v_2\), and
\(\ext{\bm y}\in\ExtensionField^{\log_2m}\)\;

\tcp{Output the column weights and tensor row factors}
Set
\(\ext{\bm\xi}_1\leftarrow\tensorspan{\ext{\bm y}}\) and
\(\ext{\bm\xi}_2\leftarrow
\ext{\bm l}\circ\tensorspan{\ext{\bm y}}\)\;
Represent the \(j\)-th factor of \(\ext{\bm\xi}_2\) by
\((\ext l_{j,0}(1-\ext y_j),\ext l_{j,1}\ext y_j)^\top\)\;
Output
\(\ext v_1\EqOrNot\ext{\bm\xi}_1^\top\bm A\ext{\bm\rho}\)
and
\(\ext v_2\EqOrNot\ext{\bm\xi}_2^\top\bm B\ext{\bm\rho}\)\;
Set
\(\mathsf{ok}\leftarrow\mathsf{ok}_0\wedge\mathsf{ok}_{\mathsf{IP}}\)\;
\Return{\((\mathsf{ok},\ext v_1,\ext v_2,\ext{\bm\rho},
\ext{\bm\xi}_1,\ext{\bm\xi}_2)\)}\;
\end{algorithm}

% !TEX root = ../main.tex
\begin{algorithm}[!t]
\linespread{0.97}\selectfont
\caption{Packed permutation proof}
\label{alg:permutation}
\KwIn{
\(N=m\kappa\), where \(m,\kappa\in2^{\mathbb N}\);
\(\bm a,\bm b\in\Field^N\)
}
\KwOut{
Terminal values \(\ext{v}_a,\ext{v}_b\) and MLE points
\(\ext{\bm y}_a,\ext{\bm y}_b
\in\ExtensionField^{\log_2N}\);
acceptance bit \(\mathsf{ok}\)
}

\tcp{Reduce multiset equality to equality of grand products}
Fix the commitments to \(\bm a\) and \(\bm b\), then sample
\(\VtoP\):
\(\ext z\ChooseRandom\ExtensionField\)\;

\ForEach{\(\bm d\in\{\bm a,\bm b\}\)}{
    Arrange \(\bm d\) columnwise as
    \(
        \bm D^{(\bm d)}
        \leftarrow
        \concat{
            \bm d^{(1)},\ldots,\bm d^{(\kappa)}
        }
        \in\Field^{m\times\kappa}.
    \)

    \(\Prover\) initializes the column-product tree with
    \(
        \bm C^{(\bm d)}_0(Z)
        \leftarrow
        Z\bm 1_{m\times\kappa}
        +
        \bm D^{(\bm d)}
    \)\;

    \For{\(t=0,\ldots,\log_2\kappa-1\)}{
        \(
            \bm C^{(\bm d)}_{t+1}(Z)
            \leftarrow
            \bm C^{(\bm d)}_{t,L}(Z)
            \circ
            \bm C^{(\bm d)}_{t,R}(Z).
        \)
    }

    \(\Prover\) evaluates the last column layer:
    \(
        \ext{\bm c}^{(\bm d)}_0
        \leftarrow
        \bm C^{(\bm d)}_{\log_2\kappa}(\ext z)
        \in\ExtensionField^m.
    \)

    \For{\(i=0,\ldots,\log_2m-1\)}{
        \(
            \ext{\bm c}^{(\bm d)}_{i+1}
            \leftarrow
            \ext{\bm c}^{(\bm d)}_{i,L}
            \circ
            \ext{\bm c}^{(\bm d)}_{i,R}.
        \)
    }

    Set
    \(
        \Theta_{\bm d}
        \leftarrow
        \ext c^{(\bm d)}_{\log_2m}.
    \)
}

\(\PtoV\):
\(\Theta_{\bm a},\Theta_{\bm b}\)\;

\(\Verifier\) checks
\(
    \mathsf{ok}_{\mathsf{perm}}
    \leftarrow
    \left(
        \Theta_{\bm a}
        \EqOrNot
        \Theta_{\bm b}
    \right).
\)

\ForEach{\(\bm d\in\{\bm a,\bm b\}\)}{
    Set
    \(
        \ext{s}^{(\bm d)}_{\log_2m}
        \leftarrow
        \Theta_{\bm d}
    \)
    and
    \(
        \ext{\bm\xi}^{(\bm d)}_{\log_2m}
        \leftarrow(1).
    \)

    \tcp{Unfold the row-product tree}
    \For{\(i=\log_2m-1,\ldots,0\)}{
        Write the current claim as
        \(
            \ext{s}^{(\bm d)}_{i+1}
            \EqOrNot
            \left\langle
                \ext{\bm c}^{(\bm d)}_{i,L},
                \ext{\bm c}^{(\bm d)}_{i,R}
                \circ
                \ext{\bm\xi}^{(\bm d)}_{i+1}
            \right\rangle.
        \)

        Apply Fact~\ref{fact:ip} and randomly combine the two
        terminal claims to obtain
        \(
            \ext{s}^{(\bm d)}_i
            \EqOrNot
            \left\langle
                \ext{\bm\xi}^{(\bm d)}_i,
                \ext{\bm c}^{(\bm d)}_i
            \right\rangle
        \)
        and the acceptance bit
        \(\mathsf{ok}^{(\bm d)}_{\mathsf{row},i}\), retaining
        the tensor factors of \(\ext{\bm\xi}^{(\bm d)}_i\)\;
    }

    Set
    \(
        \ext{q}^{(\bm d)}_{\log_2\kappa}
        \leftarrow
        \ext{s}^{(\bm d)}_0
    \)
    and
    \(
        \ext{\bm\zeta}^{(\bm d)}_{\log_2\kappa}
        \leftarrow
        \ext{\bm\xi}^{(\bm d)}_0.
    \)

    \tcp{Unfold the column-product tree}
    \For{\(t=\log_2\kappa-1,\ldots,0\)}{
        Write the current claim as
        \(
            \ext{q}^{(\bm d)}_{t+1}
            \EqOrNot
            \left\langle
                \ext{\bm\zeta}^{(\bm d)}_{t+1},
                \operatorname{vec}
                \bigl(
                    \bm C^{(\bm d)}_{t,L}(\ext z)
                \bigr)
                \circ
                \operatorname{vec}
                \bigl(
                    \bm C^{(\bm d)}_{t,R}(\ext z)
                \bigr)
            \right\rangle.
        \)

        Apply the coefficient-packing reduction of
        Algorithm~\ref{alg:hadamard} and use Fact~\ref{fact:ip}
        for its evaluated inner product to obtain
        \(
            \ext{q}^{(\bm d)}_t
            \EqOrNot
            \left\langle
                \ext{\bm\zeta}^{(\bm d)}_t,
                \operatorname{vec}
                \bigl(
                    \bm C^{(\bm d)}_t(\ext z)
                \bigr)
            \right\rangle
        \)
        and the acceptance bit
        \(\mathsf{ok}^{(\bm d)}_{\mathsf{col},t}\), retaining
        the tensor factors of \(\ext{\bm\zeta}^{(\bm d)}_t\)\;
    }

    Write
    \(
        \ext{\bm\zeta}^{(\bm d)}_0
        =
        \tensorspan{\ext{\bm y}_{\bm d}}
    \)
    for
    \(
        \ext{\bm y}_{\bm d}
        \in\ExtensionField^{\log_2N}.
    \)

    \(\Prover\) and \(\Verifier\) remove the public shift:
    \(
        \ext{v}_{\bm d}
        \leftarrow
        \ext{q}^{(\bm d)}_0
        -
        \ext z
        \left\langle
            \ext{\bm\zeta}^{(\bm d)}_0,
            \bm 1_N
        \right\rangle.
    \)

    Output the terminal claim
    \(
        \ext{v}_{\bm d}
        \EqOrNot
        \left\langle
            \bm d,
            \tensorspan{\ext{\bm y}_{\bm d}}
        \right\rangle.
    \)
}

Set
\(
    \ext{v}_a\leftarrow\ext{v}_{\bm a}
\),
\(
    \ext{v}_b\leftarrow\ext{v}_{\bm b}
\),
\(
    \ext{\bm y}_a\leftarrow\ext{\bm y}_{\bm a}
\),
and
\(
    \ext{\bm y}_b\leftarrow\ext{\bm y}_{\bm b}.
\)

Set
\(
    \mathsf{ok}
    \leftarrow
    \mathsf{ok}_{\mathsf{perm}}
    \wedge
    \bigwedge_{\bm d\in\{\bm a,\bm b\}}
    \left(
        \bigwedge_{i=0}^{\log_2m-1}
        \mathsf{ok}^{(\bm d)}_{\mathsf{row},i}
    \right)
    \wedge
    \bigwedge_{\bm d\in\{\bm a,\bm b\}}
    \left(
        \bigwedge_{t=0}^{\log_2\kappa-1}
        \mathsf{ok}^{(\bm d)}_{\mathsf{col},t}
    \right).
\)

\Return
\(
    (
        \mathsf{ok},
        \ext{v}_a,
        \ext{v}_b,
        \ext{\bm y}_a,
        \ext{\bm y}_b
    )
\)\;
\end{algorithm}

% !TEX root = ../main.tex
\begin{algorithm}[!t]
\caption{Packed opening protocol}
\label{alg:packed-pcs}
\KwPubIn{
\(\mathsf{pp}_{\pc},\mathsf{cm}\);
\(\Field\subset\ExtensionField\) with basis \(\ext{\bm\beta}\)
and \(\ext\beta_1=1\);
\(N=m\kappa\), \(\ell=\log_2m\);
\(\ext v\in\ExtensionField\),
\(\ext{\bm\rho}\in\ExtensionField^\kappa\), and
\(\ext{\bm\xi}=\bigotimes_{j=1}^{\ell}
\ext{\bm\xi}^{(j)}\), given by factors
\(\ext{\bm\xi}^{(j)}=
(\ext\xi_0^{(j)},\ext\xi_1^{(j)})^\top\in\ExtensionField^2\)
}
\KwPriIn{
\(\Prover\): \(\bm A\in\Field^{m\times\kappa}\) and
commitment randomness \(\rho_{\pc}\)
}
\KwOut{Acceptance bit \(\mathsf{ok}\)}

\(\Verifier\) returns \(0\) on a malformed message or a failed
check, reduction, or opening\;
\tcp{Authenticate the column evaluations by an extension-field projection}
\(\Prover\) sets
\(\ext{\bm a}\leftarrow\bm A\ext{\bm\beta}\) and decodes
\(\bm\Xi\in\Field^{m\times\kappa}\) from
\(\ext{\bm\xi}=\bm\Xi\ext{\bm\beta}\)\;
\(\Prover\) computes \(\bm U\leftarrow\bm\Xi^\top\bm A\)\;
\(\PtoV\): \(\ext{\bm u}\leftarrow\bm U^\top\ext{\bm\beta}
\in\ExtensionField^\kappa\)\;
\(\Verifier\) checks
\(\ext v\EqOrNot\ext{\bm u}^{\top}\ext{\bm\rho}\)\;
\(\Verifier\) recovers \(\bm U\in\Field^{\kappa\times\kappa}\)
from \(\ext{\bm u}^{\top}=\ext{\bm\beta}^{\top}\bm U\)\;
\(\VtoP\): \(\ext{\bm r}\ChooseRandom\ExtensionField^\kappa\)\;
\(\Verifier\) computes
\(\ext c\leftarrow\ext{\bm r}^{\top}\bm U\ext{\bm\beta}\)\;
\(\Prover\) computes the operand
\(\bm\Xi\ext{\bm r}\in\ExtensionField^m\)\;
\(\Prover,\Verifier\) apply Fact~\ref{fact:ip} over
\(\ExtensionField\) to
\(\ext c\EqOrNot\langle\ext{\bm a},\bm\Xi\ext{\bm r}\rangle\),
and obtain terminal values \(\ext w_{\mathsf{com}},\ext w\)
and point \(\ext{\bm y}\in\ExtensionField^\ell\)\;
\tcp{The verifier contracts the public queries directly}
Use the multiplication matrices \(\bm M_{\ext z}\) from
Eq.~\eqref{eq:pcs-multiplication-matrix}, applied by extension-field
scaling of rows and columns without materializing \(\bm M_{\ext z}\)\;
\(\Verifier\) initializes
\(\bm W\leftarrow\bm e_1\bm e_1^\top\in\Field^{\kappa\times\kappa}\),
where \(\bm e_1=(1,0,\ldots,0)^\top\)\;
\For{\(j=\ell,\ldots,1\)}{
    \(\Verifier\) updates
    \(\bm W\leftarrow
       \bm W\bm M_{\ext\xi_0^{(j)}}^\top
       +\bm M_{\ext y_j}\bm W
        \bm M_{\ext\xi_1^{(j)}-\ext\xi_0^{(j)}}^\top\)\;
    \tcp{All occurrences on the right use the previous \(\bm W\), as in Eq.~\eqref{eq:pcs-query-recurrence}}
}
\(\Verifier\) checks
\(\ext w\EqOrNot\ext{\bm\beta}^{\top}\bm W\ext{\bm r}\)\;

\(\mathsf{ok}\leftarrow
\pc.\Open\!\left(
    \mathsf{pp}_{\pc},\mathsf{cm},\ext w_{\mathsf{com}},
    \ext{\bm y}:\ext{\bm a},\rho_{\pc}
\right)\)\;
\Return \(\mathsf{ok}\)\;
\end{algorithm}

% !TEX root = ../main.tex
\begin{algorithm}[!t]
\footnotesize
\caption{PLONK proof with single-instance packing}
\label{alg:single-instance-packing}
\KwPubIn{
Parameters of Theorem~\ref{theorem:single-instance-packing};
\(\bm{\mathsf{io}}\); the preprocessed public commitment of
logical length \(m_{\mathsf{pub}}\)
}
\KwPriIn{
\(\Prover\): \(\bm w_L,\bm w_R,\bm w_O\in\Field^N\);
circuit tables and their opening data
}
\KwOut{Acceptance bit \(\mathsf{ok}\)}

\(\Verifier\) returns \(0\) on a malformed message or a failed
check, reduction, or opening\;
\tcp{Commit before sampling challenges}
Arrange \(\bm w_\ell\) as \(\bm W_\ell\in\Field^{m\times\kappa}\)
for \(\ell\in\{L,R,O\}\)\;
\(\PtoV\): one commitment to the concatenation of
\(\bm W_L\ext{\bm\beta},\bm W_R\ext{\bm\beta},
\bm W_O\ext{\bm\beta}\), of logical length \(3m\)\;
Set the virtual vector \(\bm w_M\leftarrow\bm w_L\circ\bm w_R\)\;

\tcp{Reduce the gate claims; defer all terminal openings}
\(\VtoP\): \(\ext{\bm x}\ChooseRandom\ExtensionField^{\log_2N}\);
set \(\ext{\bm\xi}\leftarrow\tensorspan{\ext{\bm x}}\)\;
\(\PtoV\): \(\ext c_L,\ext c_R,\ext c_M,\ext c_O,\ext c_C\)
as in Eq.~\eqref{eq:plonkWeightedHadamard1} and
Eq.~\eqref{eq:plonkWeightedHadamard2}\;
\(\Verifier\) checks
\(\ext c_L+\ext c_R+\ext c_M+\ext c_O+\ext c_C\EqOrNot0\)\;
\ForEach{\(\ell\in\{L,R,O\}\)}{
    Apply Algorithm~\ref{alg:hadamard} to
    \(\ext c_\ell\EqOrNot
    \langle\ext{\bm\xi},\bm q_\ell\circ\bm w_\ell\rangle\);
    retain both terminal queries\;
}
Apply Algorithm~\ref{alg:hadamard} to
\(\ext c_M\EqOrNot
\langle\ext{\bm\xi},\bm q_M\circ\bm w_M\rangle\);
retain its selector query\;
Express its product-vector query as
\(\ext v\EqOrNot
\ext{\bm l}^{\top}(\bm W_L\circ\bm W_R)\ext{\bm r}\).
Apply Algorithm~\ref{alg:hadamard} again and retain both terminals\;
Retain \(\ext c_C\EqOrNot\langle\ext{\bm\xi},\bm q_C\rangle\)\;

\tcp{Reduce the copy claims over \(\mathbb K\)}
Form \(\bm w\in\Field^L\), including the public I/O, and
\(\bm a\leftarrow\bm\varsigma+\theta_d\bm w\),
\(\bm b\leftarrow\bm\iota+\theta_d\bm w\)\;
Pad both operands with zeros to length \(4N\).
Apply Algorithm~\ref{alg:permutation} with base field \(\mathbb K\),
extension \(\ExtensionField\), and packing width \(\kappa_{\mathbb K}\),
sampling its fresh challenge
\(\ext z\ChooseRandom\ExtensionField\setminus\{0\}\)\;
Expand its two weighted terminals as in
Appendix~\ref{appx:proof-single-instance}.
Restrict to the wire segments and regroup the tensor factors
to packing width \(\kappa\)\;
\(\PtoV\): six wire values and
\(3\lceil\log_{\underline q}L\rceil\) permuted-digit values;
retain their weighted queries\;
\(\Verifier\) evaluates the identity-label, public-I/O, and padding
contributions directly and checks both recombinations\;

\tcp{Batch the witness and public openings}
For every retained query, execute the column-message and scalar-check
steps of Section~\ref{sec:pcs}.
Fix all column messages before sampling independent projection vectors\;
Project each column identity, extend its query by zeros outside
the committed block, and sum the inner-product claims in each group\;
For each group, apply Fact~\ref{fact:ip} and check its public
query-side evaluation as in Appendix~\ref{appx:proof-single-instance}\;
Authenticate the two terminal values by underlying PCS openings
at logical lengths \(3m\) and \(m_{\mathsf{pub}}\)\;
\Return \(1\)\;
\end{algorithm}

% !TEX root = ../main.tex
\begin{algorithm}[!t]
\footnotesize
\linespread{0.97}\selectfont
\caption{PLONK proof with packing across instances}
\label{alg:multiple-instance-packing}
\KwPubIn{
Parameters of Theorem~\ref{theorem:multiple-instance-packing};
the \(\kappa\) public-I/O vectors;
the preprocessed commitment to \(\bm Q\ext{\bm\beta}\in\ExtensionField^N\)
}
\KwPriIn{
\(\Prover\): \(\bm W_L,\bm W_R,\bm W_O\in\Field^{N\times\kappa}\);
circuit tables and their opening data
}
\KwOut{Acceptance bit \(\mathsf{ok}\)}

\(\Verifier\) returns \(0\) on a malformed message or a failed
check, reduction, or opening\;
\tcp{Commit before sampling challenges}
Set \(\ext{\bm w}_\ell\leftarrow\bm W_\ell\ext{\bm\beta}\)
for \(\ell\in\{L,R,O\}\), and the virtual vector
\(\ext{\bm w}_M\leftarrow(\bm W_L\circ\bm W_R)\ext{\bm\beta}\)\;
\(\PtoV\): one commitment to the concatenation of
\(\ext{\bm w}_L,\ext{\bm w}_R,\ext{\bm w}_O\),
of logical length \(3N\)\;

\tcp{Reduce the gate claims; defer all terminal openings}
\(\VtoP\): \(\ext{\bm x}\ChooseRandom\ExtensionField^{\log_2N}\)\;
Run sumcheck with initial claim \(0\) and individual degree at most
three for the \(\tensorspan{\ext{\bm x}}\)-weighted residual in
Eq.~\eqref{eq:together-multiple-gates};
obtain \(\ext{\bm y},\ext c\)\;
\(\PtoV\): selector values \(\ext v_\ell\) for
\(\ell\in\{L,R,M,O,C\}\) and wire values \(\ext u_\ell\) for
\(\ell\in\{L,R,M,O\}\), at \(\ext{\bm y}\)\;
\(\Verifier\) checks
\(\ext c\EqOrNot
\langle\tensorspan{\ext{\bm x}},\tensorspan{\ext{\bm y}}\rangle
[\sum_{\ell\in\{L,R,M,O\}}\ext v_\ell\ext u_\ell+
(\sum_{j=1}^{\kappa}\ext\beta_j)\ext v_C]\)\;
Retain the five selector queries and three ordinary wire queries\;
Apply Algorithm~\ref{alg:hadamard} to
\(\ext u_M\EqOrNot
\tensorspan{\ext{\bm y}}^\top
(\bm W_L\circ\bm W_R)\ext{\bm\beta}\);
retain its two weighted terminal queries\;

\tcp{Reduce the copy claims by two product trees}
Form \(\ext{\bm w}\in\ExtensionField^L\), including the packed public I/O\;
\(\VtoP\): fresh independent
\(\ext y,\ext z\ChooseRandom\ExtensionField\)\;
Construct product trees on
\(\ext z\bm1_L+\bm\iota+\ext y\ext{\bm w}\) and
\(\ext z\bm1_L+\bm\varsigma+\ext y\ext{\bm w}\),
padding each input with ones to length \(4N\)\;
\(\PtoV\): a common root \(\ext\Theta\)\;
For each tree, run the degree-three layer sumchecks of
Appendix~\ref{appx:proof-multiple-instance}, starting at \(\ext\Theta\).
After checking the sumcheck terminal identity and fixing both child values,
sample a fresh merge
challenge and continue to the next layer; retain the terminal input claim\;
Expand both terminals, retaining the prefix weights of the three
wire segments. \(\PtoV\): six wire values and three permuted-label values\;
\(\Verifier\) evaluates the identity-label, public-I/O, and padding
contributions directly and checks both recombinations\;

\tcp{Batch the witness and public openings}
Express the five selector and three permuted-label queries as weighted
queries to \(\bm Q\), as in Appendix~\ref{appx:proof-multiple-instance}\;
Execute the column-message and scalar-check steps of Section~\ref{sec:pcs}
for these eight queries and the two Hadamard terminals\;
Fix all column messages and the nine ordinary wire values.
Sample independent projection vectors for the packed queries and
independent scalar weights for the ordinary queries\;
Project the column identities and combine all inner-product claims
within the witness and public groups, extending queries by zeros
outside their committed blocks\;
For each group, apply Fact~\ref{fact:ip} and verify the combined query-side
evaluation, retaining all witness block-prefix weights\;
Authenticate the two terminal values by underlying PCS openings
at logical lengths \(3N\) and \(N\)\;
\Return \(1\)\;
\end{algorithm}


\clearpage

\subsection{Packed Polynomial Commitment Scheme}
\label{appx:pcs}

% !TEX root = main.tex
We express the ring-switching construction of Diamond and
Posen~\cite{eurocrypt/DiamondP26} in our vector and matrix notation
to support the weighted tensor queries of Section~\ref{sec:blocks}.

Let \(N=m\kappa\) and \([\ExtensionField:\Field]=\kappa\).
We commit to the fixed-basis packing \(\bm A\ext{\bm\beta}\)
of \(\bm A\in\Field^{m\times\kappa}\) and reduce each query
to one MLE opening. The verifier evaluates the resulting public
query by a recursion on small base-field matrices.
We then instantiate the underlying PCS from a base-field linear
code via Blaze's code-switching construction.

\subsubsection{Given a PCS over the Extension Field}

\medskip\noindent\textbf{Commitment and opening relation.}

Let \(m,\kappa\) be powers of two, with \(\kappa\geq2\), and write
\(\ell=\log_2m\).
Fix an \(\Field\)-basis
\(\ext{\bm\beta}=(\ext\beta_1,\ldots,\ext\beta_\kappa)^\top\)
of \(\ExtensionField\), with \(\ext\beta_1=1\), and let \(\pc\)
be a multilinear PCS for vectors in \(\ExtensionField^m\).
Extension-field elements are stored in this basis, so packing
and coordinate extraction only reorganize their base-field coordinates.
The schoolbook cost bounds below use a polynomial basis.

\paragraph{Setup and commitment.}
Run \(\pc.\Setup(1^\lambda,m)\) to obtain
\(\mathsf{pp}_{\pc}\), and include the fields, basis, and dimensions
\(N,m,\kappa\) in the public parameters.
For the witness matrix
\(\bm A=\concat{\bm a^{(1)},\ldots,\bm a^{(\kappa)}}
\in\Field^{m\times\kappa}\), pack each row into one
extension-field element, obtaining
\(\ext{\bm a}=\bm A\ext{\bm\beta}\in\ExtensionField^m\).
Compute the commitment
\(\mathsf{cm}=\pc.\Com(\mathsf{pp}_{\pc},\ext{\bm a};
\rho_{\pc})\), where \(\rho_{\pc}\) is the commitment
randomness, and retain \(\bm A\) and \(\rho_{\pc}\).
The fixed basis makes binding to \(\ext{\bm a}\) equivalent
to binding to \(\bm A\).

\paragraph{Opening relation.}
A query specifies column weights
\(\ext{\bm\rho}\in\ExtensionField^\kappa\), a tensor vector
\(\ext{\bm\xi}\in\ExtensionField^m\), and a claimed value
\(\ext v\in\ExtensionField\). The opening relation is:
\begin{equation}
    \ext v
    =\langle\bm A\ext{\bm\rho},\ext{\bm\xi}\rangle
    =\ext{\bm\xi}^{\top}\bm A\ext{\bm\rho}.
\label{eq:pcs-opening-relation}
\end{equation}
The row query is supplied as
\(\ext{\bm\xi}=\bigotimes_{j=1}^{\ell}\ext{\bm\xi}^{(j)}\),
where each factor is
\(\ext{\bm\xi}^{(j)}=(\ext\xi^{(j)}_0,
\ext\xi^{(j)}_1)^\top\in\ExtensionField^2\).
Ordinary MLE queries use the tensor spans of the corresponding
parts of the evaluation point for both \(\ext{\bm\rho}\) and
\(\ext{\bm\xi}\). Allowing arbitrary two-entry factors also
supports the coordinatewise products of tensor queries arising
in the weighted Hadamard reduction.

\medskip\noindent\textbf{Packed opening protocol.}

The commitment uses the fixed weights \(\ext{\bm\beta}\), whereas
the query uses \(\ext{\bm\rho}\). The prover first sends the
column evaluations. Coordinate decomposition and a random
projection then reduce their correctness to an inner-product
claim about the committed vector.

\paragraph{Reduce the column claims.}
Write the public query and the claimed column values in base-field
coordinates:
\[
    \ext{\bm\xi}=\bm\Xi\ext{\bm\beta},
    \qquad
    \ext{\bm u}^{\top}=\ext{\bm\beta}^{\top}\bm U,
    \qquad
    \bm\Xi\in\Field^{m\times\kappa},
    \quad
    \bm U\in\Field^{\kappa\times\kappa}.
\]
The honest prover computes \(\bm U=\bm\Xi^\top\bm A\) and sends
\(\ext{\bm u}=\bm U^\top\ext{\bm\beta}\).
The verifier extracts \(\bm U\) from this message and checks
\(\ext v=\ext{\bm u}^{\top}\ext{\bm\rho}\).
The matrix \(\bm\Xi\) is defined implicitly by the public query.

Since \(\bm A\) has base-field entries, the remaining claim
\(\ext{\bm u}^{\top}=\ext{\bm\xi}^{\top}\bm A\) is equivalent
to \(\bm U=\bm\Xi^\top\bm A\). Packing the rows of this
identity introduces the committed vector \(\ext{\bm a}\):
\begin{equation}
    \bm U=\bm\Xi^\top\bm A
    \quad\Longleftrightarrow\quad
    \bm U\ext{\bm\beta}
    =\bm\Xi^\top\ext{\bm a}.
\label{eq:pcs-encoded-matrix}
\end{equation}
After receiving \(\ext{\bm u}\), the verifier samples
\(\ext{\bm r}\sample\ExtensionField^\kappa\), independently of
the original query weights, and projects
Eq.~\eqref{eq:pcs-encoded-matrix} to one inner-product claim:
\begin{equation}
    \ext c
    =\ext{\bm r}^{\top}\bm U\ext{\bm\beta}
    =\langle\ext{\bm a},\bm\Xi\ext{\bm r}\rangle.
\label{eq:pcs-random-opening}
\end{equation}
The verifier computes \(\ext c\) by packing the rows of
\(\bm U\) and taking their inner product with \(\ext{\bm r}\).

Apply Fact~\ref{fact:ip} over \(\ExtensionField\) to
Eq.~\eqref{eq:pcs-random-opening}.
Let \(\ext{\bm y}\in\ExtensionField^\ell\) be its terminal point.
The prover's terminal values
\(\ext w_{\mathsf{com}},\ext w\in\ExtensionField\) give two claims:
\begin{equation}
    \ext w_{\mathsf{com}}
    =\langle\ext{\bm a},\tensorspan{\ext{\bm y}}\rangle,
    \qquad
    \ext w=\tensorspan{\ext{\bm y}}^{\top}\bm\Xi\ext{\bm r}.
\label{eq:pcs-terminal-claims}
\end{equation}
The first is an MLE opening of the committed vector.
The verifier checks the second by the following matrix recursion.

\paragraph{Evaluate the public query.}
Write \(\tensorspan{\ext{\bm y}}=\bm Y\ext{\bm\beta}\), where
\(\bm Y\in\Field^{m\times\kappa}\).
It suffices to compute the small matrix:
\begin{equation}
    \bm W=\bm Y^\top\bm\Xi\in\Field^{\kappa\times\kappa},
    \qquad
    \ext w=\ext{\bm\beta}^{\top}\bm W\ext{\bm r}.
\label{eq:pcs-query-matrix}
\end{equation}
The public tensor factors determine \(\bm Y\) and \(\bm\Xi\),
so the verifier can compute their product without forming either
length-\(m\) coordinate table.

For \(\ext z\in\ExtensionField\), let
\(\bm M_{\ext z}\in\Field^{\kappa\times\kappa}\) be the coordinate
matrix of multiplication by \(\ext z\), characterized by:
\begin{equation}
    \ext{\bm\beta}^{\top}\bm M_{\ext z}
    =\ext z\,\ext{\bm\beta}^{\top}.
\label{eq:pcs-multiplication-matrix}
\end{equation}
Thus left multiplication by \(\bm M_{\ext z}\) scales each
coordinate column by \(\ext z\) in \(\ExtensionField\), and right
multiplication by \(\bm M_{\ext z}^\top\) scales each coordinate row.
In particular, \(\bm M_{1-\ext z}=\bm I-\bm M_{\ext z}\).

Let \(\bm Y_j,\bm\Xi_j\) be the coordinate matrices for the
suffix factors \(j,\ldots,\ell\):
\[
    \bm Y_j\ext{\bm\beta}
    =\bigotimes_{i=j}^{\ell}
        \begin{pmatrix}1-\ext y_i\\ \ext y_i\end{pmatrix},
    \qquad
    \bm\Xi_j\ext{\bm\beta}
    =\bigotimes_{i=j}^{\ell}\ext{\bm\xi}^{(i)}.
\]
Here \(\bm Y_1=\bm Y\), \(\bm\Xi_1=\bm\Xi\), and
\(\bm Y_{\ell+1}=\bm\Xi_{\ell+1}=\bm e_1^\top\), where
\(\bm e_1=(1,0,\ldots,0)^\top\in\Field^\kappa\).
The two possible values of the first suffix bit give the block
decompositions:
\begin{equation}
\begin{aligned}
    \bm Y_j
    &=
    \begin{bmatrix}
        \bm Y_{j+1}\bm M_{1-\ext y_j}^\top\\
        \bm Y_{j+1}\bm M_{\ext y_j}^\top
    \end{bmatrix},\\
    \bm\Xi_j
    &=
    \begin{bmatrix}
        \bm\Xi_{j+1}\bm M_{\ext\xi^{(j)}_0}^\top\\
        \bm\Xi_{j+1}\bm M_{\ext\xi^{(j)}_1}^\top
    \end{bmatrix}.
\end{aligned}
\label{eq:pcs-query-blocks}
\end{equation}
Consequently, the verifier computes
\(\bm W_j=\bm Y_j^\top\bm\Xi_j\) by initializing
\(\bm W_{\ell+1}=\bm e_1\bm e_1^\top\) and processing
\(j=\ell,\ldots,1\):
\begin{equation}
\begin{aligned}
    \bm W_j
    &=
    \bm M_{1-\ext y_j}\bm W_{j+1}
        \bm M_{\ext\xi^{(j)}_0}^\top
    +\bm M_{\ext y_j}\bm W_{j+1}
        \bm M_{\ext\xi^{(j)}_1}^\top\\
    &=
    \bm W_{j+1}\bm M_{\ext\xi^{(j)}_0}^\top
    +\bm M_{\ext y_j}\bm W_{j+1}
        \bm M_{\ext\xi^{(j)}_1-\ext\xi^{(j)}_0}^\top.
\end{aligned}
\label{eq:pcs-query-recurrence}
\end{equation}
This recursion yields \(\bm W=\bm W_1\) without forming
the suffix coordinate matrices.
Each step of the second line uses at most \(3\kappa\)
extension-field multiplications, implemented by row and column
scaling without materializing \(\bm M_{\ext z}\).
The recursion takes \(O(\kappa\ell)\) extension-field
multiplications and \(O(\kappa^2\ell)\) base-field additions and
coordinate moves, using \(O(\kappa^2)\) base-field coordinates
of working memory. It applies in every characteristic.

The verifier packs the columns of \(\bm W\), computes
\(\ext{\bm\beta}^{\top}\bm W\ext{\bm r}\), and checks the second
identity in Eq.~\eqref{eq:pcs-query-matrix}.
This computation is exact and requires no additional prover message.

\paragraph{Terminal opening.}
Finally, the parties authenticate \(\ext w_{\mathsf{com}}\)
at \(\ext{\bm y}\) by invoking
\(\pc.\Open(\mathsf{pp}_{\pc},\mathsf{cm},
\ext w_{\mathsf{com}},\ext{\bm y}
:\ext{\bm a},\rho_{\pc})\).
The verifier accepts if the scalar checks, the inner-product
reduction, and the underlying opening accept.
Algorithm~\ref{alg:packed-pcs} in Appendix~\ref{appx:pseudocode}
summarizes the protocol.

\medskip\noindent\textbf{Security and complexity.}

\begin{proposition}[Packed polynomial commitment scheme]
\label{proposition:packed-pcs}
Let \(\pc\) bind the prover to a vector in \(\ExtensionField^m\),
with evaluation-soundness error \(\epsilon_{\mathsf{PCS}}\) for
adaptively chosen MLE claims about that vector.
The packed commitment is binding and the opening protocol is
complete, with evaluation-soundness error at most:
\begin{equation}
    \epsilon_{\mathsf{PCS}}
    +\frac{1+2\ell}{|\ExtensionField|}.
\label{eq:pcs-soundness}
\end{equation}
The opening reduction requires no additional lower bound on
\(|\Field|\).

Excluding the underlying PCS opening, the prover's arithmetic cost is:
\begin{equation}
    O(\kappa N)\,\bb
    +O(N)\,\be
    +O(m+\kappa)\,\ee.
\label{eq:pcs-prover-cost}
\end{equation}
The verifier's arithmetic cost is:
\begin{equation}
    O\bigl(\kappa^2(1+\ell)\bigr)\,\bb
    +O\bigl(\kappa(1+\ell)\bigr)\,\ee.
\label{eq:pcs-verifier-cost}
\end{equation}
The verifier bound includes base-field additions and coordinate
reorganization in the matrix recursion.
When \(m\geq\kappa\), schoolbook arithmetic gives
\(O(\kappa N)\) base-field operations for the prover.
Peak prover memory is \(O(N+\kappa^2)\) base-field coordinates,
or \(O(N)\) when \(m\geq\kappa\), excluding the workspace of the
underlying PCS.

With \(\ext v\) already part of the public statement,
prover-to-verifier communication, excluding the underlying
opening, consists of \(\kappa+2\ell+2\) extension-field elements.
The packed commitment inherits hiding from the underlying
commitment. The opening protocol reveals the column evaluations
\(\ext{\bm u}\), so this inheritance does not establish
zero knowledge of the opening transcript.
\end{proposition}

The proof is deferred to Appendix~\ref{appx:proof-pcs}.

\subsubsection{Given a Linear Code over the Base Field}

A base-field generator matrix commutes with basis packing and
therefore also encodes packed messages over \(\ExtensionField\).
We instantiate the preceding construction via Blaze's code-switching
protocol~\cite{eurocrypt/BrehmCFRRZ25}, with a BaseFold
code~\cite{crypto/ZeilbergerCF24} as the target, under Blaze's
source-encoding and proximity assumptions.

Let \(\Code:\Field^m\to\Field^n\) be a linear code with
generator matrix \(\bm G\in\Field^{n\times m}\).
Its \(\kappa\)-way interleaving encodes the columns of
\(\bm A\in\Field^{m\times\kappa}\) independently.
Packing the resulting codeword in the fixed basis gives:
\begin{equation}
\begin{aligned}
    \bm C&=\bm G\bm A\in\Field^{n\times\kappa},\\
    \ext{\bm c}&=\bm C\ext{\bm\beta}
       =\bm G(\bm A\ext{\bm\beta})
       =\Code_{\ExtensionField}(\ext{\bm a}).
\end{aligned}
\label{eq:pcs-code-packing}
\end{equation}
Here \(\Code_{\ExtensionField}\) is the scalar extension defined
by the same generator matrix; it has the same rate and relative
distance as \(\Code\).
Thus a row of \(\bm C\) represents one extension-field codeword
symbol. Committing to these rows supplies Blaze's source codeword
oracle for \(\ext{\bm a}=\bm A\ext{\bm\beta}\).

Apply Blaze's code-switching construction over \(\ExtensionField\)
to this oracle, switching to a BaseFold encoding of the same
length-\(m\) message \(\ext{\bm a}\).
The resulting MLE openings instantiate the underlying PCS \(\pc\)
in the preceding construction.
The packed opening reduction handles the challenge-dependent weights
\(\ext{\bm\rho}\), and code switching authenticates its terminal
fixed-basis opening. We include the costs and soundness error of
the chosen Blaze instantiation, including source-encoding
consistency, in the underlying PCS costs and
\(\epsilon_{\mathsf{PCS}}\), respectively.

If \(\Code\) is constant-rate and linear-time encodable, forming
\(\bm C\) costs \(O(N)\) base-field operations and its payload
occupies \(O(N\log_2|\Field|)\) bits, excluding Merkle-tree hashes.
The BaseFold stage works on \(m=N/\kappa\) extension-field message
symbols. Thus the source encoding retains the witness's small-field
representation, and the extension-field opening operates on the
shorter packed message.


\clearpage

\iffalse
\subsection{Complexity Bounds for the Hadamard Comparison}
\label{appx:hadamard-comparison}

Table~\ref{tab:hadamard-comparison} compares proofs of
\(\bm a\circ\bm b=\bm c\) for length-\(N\) vectors.
The constraint \(a_i b_i-c_i\) has degree two. Costs include equality
weights, the linear output check, witness-dependent preprocessing, and all reduction
rounds, but exclude commitments and openings. The derived upper bounds
use schoolbook arithmetic: multiplication in
\(\Field\times\Field\), \(\Field\times\ExtensionField\), and
\(\ExtensionField\times\ExtensionField\) costs \(O(1)\),
\(O(\kappa)\), and \(O(\kappa^2)\), respectively, where
\(\kappa=[\ExtensionField:\Field]\). RAM counts base-field elements
including inputs and read-only tables, not measured RSS.
Public field and basis preprocessing is separate; lookup-table
construction is accounted for below. Table accesses are charged
per base-field coordinate, as in Section~\ref{sec:notation}.

\smallskip\noindent\textbf{Optimized rounds and omitted terms.}
Let \(0\leq t\leq\log_2N\) count optimized or skipped rounds.
For small-field-first~\cite{iacr/BagadDT24,iacr/BagadDDT25,iacr/DaoDBDT26},
streamed prefix binding and a weighted \(3^t\)-entry quadratic grid give:
\begin{equation}
\begin{aligned}
\text{Arithmetic:}\quad&
O\!\left(\kappa N(3/2)^t+\kappa^2N/2^t+\kappa^2 3^t\right),\\
\text{RAM:}\quad& O\!\left(N+\kappa N/2^t+\kappa 3^t\right).
\end{aligned}
\label{eq:hadamard-sff-full-cost}
\end{equation}
The last terms account for grid contractions and storage.
For \(N\geq6^t\), they are absorbed by the corresponding
\(N/2^t\) terms, yielding the table entries.
For univariate skip~\cite{iacr/Gruen24}, including FFT interpolation,
weighting, binding, first-round processing, and the sumcheck tail gives:
\begin{equation}
\begin{aligned}
\text{Arithmetic:}\quad&
O\!\left(Nt+\kappa N+\kappa^2N/2^t+\kappa^2 2^t\right),\\
\text{RAM:}\quad& O\!\left(N+\kappa N/2^t+\kappa 2^t\right).
\end{aligned}
\label{eq:hadamard-skip-full-cost}
\end{equation}
Here \(N\geq4^t\) suffices to absorb the last terms.
The \(Nt\) term uses FFT interpolation on suitable subgroup
and off-domain evaluation sets. The protocol also admits general
interpolation sets~\cite[Section~5.1, footnote~18]{iacr/Gruen24},
with their interpolation cost accounted for separately. When supported,
\(2^t=\Theta(\kappa)\) and \(N=\Omega(\kappa^2)\) give
\(O(\kappa N)\) arithmetic and \(O(N)\) RAM.
For \(t\) fixed by \(\kappa\), the omitted terms are independent of
\(N\). If the absorption conditions fail, use the full bounds in
Eq.~\eqref{eq:hadamard-sff-full-cost} and
Eq.~\eqref{eq:hadamard-skip-full-cost}.
For degree-two constraints, the off-domain first-round values in
Gruen's protocol and the quotient coefficients in our Hadamard
reduction represent the same vanishing residual. Its known zeros
on the interpolation set determine the conversion between the two
representations~\cite[Section~6]{iacr/Gruen24}.

\smallskip\noindent\textbf{Flock's lookup-based skip.}
We re-cost the degree-two zerocheck of
Flock~\cite[Sections~4.2--4.4]{iacr/BunzRW26} with a conventional
sumcheck tail. Fix its byte-sized chunks and subfield
\(\mathbb F_{256}\subseteq\ExtensionField\), and let
\(3\leq t\leq7\) be the number of skipped variables.
Each row has \(2^t\) input bits and \(2^t\) off-domain samples.
Its LDE uses \(2^t/8\) lookups of length-\(2^t\) byte vectors,
so all rows cost \(O(N2^t)\) bit operations, including reads
and XORs. Byte-field products cost \(O(N)\).
In compatible \(\ExtensionField/\mathbb F_{256}\) coordinates,
weighted aggregation costs \(O(\kappa N)\), as do binding at
the skip challenge and the linear output check.
The conventional tail and its equality weights cost
\(O(\kappa^2N/2^t)\). Including parameter-only message
processing and fixed lookup tables gives:
\begin{equation}
\begin{aligned}
\text{Arithmetic:}\quad&
O\!\left(N2^t+\kappa N+\kappa^2N/2^t+\kappa^2 2^t\right),\\
\text{RAM:}\quad&
O\!\left(N+\kappa N/2^t+\kappa 2^t\right).
\end{aligned}
\label{eq:hadamard-flock-full-cost}
\end{equation}
The last terms are absorbed when \(N\geq4^t\).
This bound covers Flock's algebraic reduction to evaluation claims.
Its implementation uses \(t=6\), \(\kappa=128\),
a 16\,KiB compressed LDE table, and a 64\,KiB table for its
fixed equality weights. Compression improves table storage and
cache locality, not the asymptotic LDE operation count.
The admissible skip length is bounded by the byte-field domain,
independently of \(\kappa\).

\smallskip\noindent\textbf{Binary-tower instantiations.}
Assume \(N\geq\kappa\). The Dao--Thaler row pairs its reduction
with conventional sumcheck~\cite{iacr/DaoT24}.
The Liu--Zhang row uses every other level of a Wiedemann
tower~\cite{iacr/LiuZ25}, with mixed subfield/top-field multiplication
and streamed deterministic weights~\cite{iacr/DaoT24}.
RAM retains folded inputs and random-suffix weight tables.

\smallskip\noindent\textbf{Packing assumptions.}
The RMFE--Aurora bound~\cite{iacr/CascudoG21} includes sparse Boolean
R1CS and Reed--Solomon encoding. It assumes constant-rate packing and
codes, linear-space quasilinear FFTs, security \(O(\kappa)\), and a
packed instance large enough to absorb security padding.
Packed sumcheck~\cite{iacr/WeiWXZDWZ25} uses total packed length \(N\),
packing/repetition parameter \(r\), and \(|\Field|>3(2r-1)\).
Streaming forward differences retains \(O(r^2)\) first-round data.
This term remains in the table because \(r\) is independent of
\(\kappa\); it is absorbed if \(N\geq r^2\).
For the base-field \methodname{} row, \(|\Field|\geq\kappa\),
\(N=m\kappa\), \(m\geq\kappa\), and precomputed
bases give \(O(\kappa N)\) arithmetic and
\(O(N+\kappa^2)=O(N)\) RAM, including the linear output check.
All parameter choices must satisfy the target soundness bound.
In particular, \(r=\Theta(\kappa)\) is not assumed, and univariate
skip's error depends on \(2^t\).

\smallskip\noindent\textbf{Intermediate-field accounting.}
Let
\(t_{\mathbb K}=[\mathbb K:\Field]\) be the degree chosen in
Section~\ref{sec:hadamard}. Without lookup tables, multiplication in
\(\Field\times\mathbb K\) costs \(O(t_{\mathbb K})\)
base-field operations, whereas multiplication in \(\mathbb K\)
costs \(O(t_{\mathbb K}^2)\).
In Eq.~\eqref{eq:hadamard-small-field-aggregation}, computing
\(\bm A\bm P\) and \(\bm B\bm P\) therefore costs
\(O(N\kappa t_{\mathbb K})\), and their entrywise product costs
\(O(Nt_{\mathbb K}^2)\).
The contraction with \(\bm L\) uses
\(O(N\kappa/t_{\mathbb K})\) multiplications in \(\mathbb K\),
or \(O(N\kappa t_{\mathbb K})\) base-field operations.
Only the \(\kappa/t_{\mathbb K}\) aggregate coordinate tables
are interpolated. Dense interpolation and division on each cost
\(O(\kappa^2)\) operations in \(\mathbb K\), for a total of
\(O(\kappa^3t_{\mathbb K})\) base-field operations.
The column claims, subsequent inner-product reduction, and linear
output check cost \(O(N\kappa+\kappa^3)\).
Thus Eq.~\eqref{eq:hadamard-small-field-cost} includes the
online parameter-only work; \(N\geq\kappa^2\) gives the
untabulated bound \(O(N\kappa\log_q\kappa)\).

For the space bound, generate one column of \(\bm P\) at a time
using prefix and suffix products of the node differences, and
scan the original rows to accumulate its column of \(\bm U\).
With the public Lagrange denominators precomputed, generating all
columns costs \(O(\kappa^2)\) operations in \(\mathbb K\),
absorbed by the bounds above. The witness and \(\bm L\) use
\(O(N)\) base-field coordinates, and \(\bm U\) uses
\(O(\kappa^2)\). Interpolation basis polynomials can likewise
be generated one at a time from the public vanishing polynomial,
so no dense interpolation matrix is stored.
Fixed conversions between the tower coordinates and the original
commitment basis take at most \(O(N\kappa+\kappa^3)\) operations
and \(O(\kappa^2)\) additional coordinates.
The resulting RAM bound is \(O(N+\kappa^2)\), with the witness
still stored over \(\Field\).

\smallskip\noindent\textbf{Tabulated intermediate-field implementation.}
The second \methodname{} row uses
Eq.~\eqref{eq:hadamard-block-tabulated-cost}.
Choose \(\tau=\lfloor\log_q\kappa\rfloor\), so
\(q^\tau\leq\kappa\) and \(t_{\mathbb K}/\tau=O(1)\).
The block tables of Eq.~\eqref{eq:hadamard-block-table}
form all samples in \(O(N\kappa t_{\mathbb K}/\tau)
=O(N\kappa)\) base-field operations.
Each multiplication-table lookup reads and writes
\(O(t_{\mathbb K})\) coordinates. Thus the
\(O(N\kappa/t_{\mathbb K})\) products in the contraction
cost \(O(N\kappa)\), and aggregate interpolation and division
cost \(O(\kappa^3)\). Samplewise products and the remaining
reduction steps are absorbed in these bounds. Hence:
\begin{equation}
\begin{aligned}
T_{\mathrm{online}}&=O(N\kappa+\kappa^3),\\
T_{\mathrm{pre}}&=O(\kappa^3+t_{\mathbb K}^2q^{2t_{\mathbb K}}),\\
S_{\mathrm{total}}&=O(N+\kappa^3+t_{\mathbb K}q^{2t_{\mathbb K}}).
\end{aligned}
\label{eq:hadamard-tabulation-full-cost}
\end{equation}
The multiplication table has \(q^{2t_{\mathbb K}}\) entries;
schoolbook construction costs \(O(t_{\mathbb K}^2)\) per entry.
The block tables use \(O(\kappa^3)\) coordinates and operations:
precompute scalar multiples of the sample-map rows, then enumerate
input blocks by extending prefixes. Both tables depend only on
public field and interpolation parameters and are reused across
witnesses. Streaming witness rows leaves \(O(N+\kappa^2)\)
working coordinates. For \(N\geq\kappa^2\), the online time is
\(O(N\kappa)\), but the read-only tables remain in the RAM entry.
\fi